\subsection{CONTINUOUS MULTILINEAR FUNCTIONS}

THEOREM 1. Let \(\mathbf{E}_i(\mathbf{i} \leqslant i \leqslant n)\) and \(\mathbf{F}\) be normed spaces over a non-discrete valued division ring \(\mathbf{K}\) , and let \(f\) be a multilinear mapping of \(\prod_{i=1}^{n} \mathbf{E}_i\) into \(\mathbf{F}\) . Then \(f\) is continuous on \(\prod_{i=1}^{n} \mathbf{E}_i\) if and only if there exists a real number \(a > 0\) such that, for all \(\boldsymbol{x}_i \in \mathbf{E}_i\) ( \(\mathbf{i} \leqslant i \leqslant n\) ), we have

\[
\left| \left| f \left(\boldsymbol {x} _ {1}, \boldsymbol {x} _ {2}, \dots , \boldsymbol {x} _ {n}\right) \right| \right| \leqslant a. \left| \left| \boldsymbol {x} _ {1} \right| \right|. \left| \left| \boldsymbol {x} _ {2} \right| \right|. \dots . \left| \left| \boldsymbol {x} _ {n} \right| \right|.\tag{3}
\]

The condition is necessary. For if \(f\) is continuous at the point \((\mathbf{o}, \mathbf{o}, \ldots, \mathbf{o})\) there exists a number \(b > 0\) such that the relations \(||\boldsymbol{x}_i|| \leqslant b\) ( \(\mathbf{i} \leqslant i \leqslant n\) ) imply \(||f(\boldsymbol{x}_1, \ldots, \boldsymbol{x}_n)|| \leqslant \mathbf{i}\) . Let \(t_0\) be an element of \(\mathbf{K}\) such that \(0 < |t_0| < \mathbf{i}\) ; then for every point \((\boldsymbol{x}_i) \in \prod_{i=1}^{n} \mathrm{E}_i\) such that none of the \(\boldsymbol{x}_i\) is zero, there exist \(n\) rational integers \(k_i\) such that \(b|t_0| < ||t_0^k |\boldsymbol{x}_i|| \leqslant b\) ; consequently we have

\[
\left| t _ {0} \right| ^ {k _ {1} + k _ {2} + \dots + k _ {n}} \left| \left| f \left(\boldsymbol {x} _ {1}, \boldsymbol {x} _ {2}, \dots , \boldsymbol {x} _ {n}\right) \right| \right| \leqslant 1;
\]

on the other hand we have \(\frac{\mathbf{I}}{|t_0|^k_i} \leqslant \frac{\mathbf{I}}{b|t_0|} ||\boldsymbol{x}_i||\) , and the relation (3) therefore follows, with \(a = (\mathbf{I}/b|t_0|)^n\) . This relation is evidently still valid when one of the \(\boldsymbol{x}_i\) is zero.

The condition is sufficient. We shall show that, if it is satisfied, \(f\) is continuous at every point \((\pmb{a}_i)\) of \(\prod_{i=1}^{n} \mathbf{E}_i\) . We can write

\[
f \left(x _ {1}, \dots , x _ {n}\right) - f \left(a _ {1}, \dots , a _ {n}\right) = \sum_ {i = 1} ^ {n} f \left(a _ {1}, \dots , a _ {i - 1}, x _ {i} - a _ {i}, x _ {i + 1}, \dots , x _ {n}\right).
\]

Now, using (3), the conditions \(||\pmb{x}_i - \pmb{a}_i|| \leqslant r\) ( \(1 \leqslant i \leqslant n\) ) imply that

\[
\left| \left| f \left(\boldsymbol {a} _ {1}, \dots , \boldsymbol {a} _ {i - 1}, \boldsymbol {x} _ {i} - \boldsymbol {a} _ {i}, \boldsymbol {x} _ {i + 1}, \dots , \boldsymbol {x} _ {n}\right) \right| \right| \leqslant a r \prod_ {k \neq i} ^ {n} \left(\left| \left| \boldsymbol {a} _ {k} \right| \right| + r\right);
\]

hence, if \(c\) is the maximum of the numbers \(\| \pmb{a}_i \| (\mathbf{i} \leqslant i \leqslant n)\) , we have

\[
\left| \left| f \left(\boldsymbol {x} _ {1}, \dots , \boldsymbol {x} _ {n}\right) - f \left(\boldsymbol {a} _ {1}, \dots , \boldsymbol {a} _ {n}\right) \right| \right| \leqslant n a r (c + r) ^ {n - 1}.
\]

Since the right-hand side of this inequality is a polynomial in r with zero constant term, it tends to o as r tends to o; hence f is continuous.

Remark. Two of the propositions proved earlier are consequences of this theorem: the continuity of the bilinear function \(tx\) , by virtue of the relation \(||tx|| = |t|.||x||\) ; and Proposition 7, by applying Theorem 1 to the identity mapping of E, considered as a linear mapping of the space E, endowed with the norm \(p\) into the space E endowed with the norm \(q\) (or vice versa).

